% C-owned proof fragment, 2026-11-19.  No historical identity is assumed.
\paragraph{Universal absolute flatness and finite piecewise morphisms.}

Let $\operatorname{Abs}_k$ be the full category of $k$-schemes whose
local rings are fields.  For a scheme $X$, let $t_X:T(X)\to X$ be its
universal absolutely flat scheme.  The functor $T$ is right adjoint to
the inclusion of absolutely flat schemes into schemes.  Thus
\[
 \operatorname{Hom}_{\operatorname{Abs}_k}(Z,T(X))
   =\operatorname{Hom}_{\operatorname{Sch}_k}(Z,X)
 \quad (Z\in\operatorname{Abs}_k).
\]
The scheme construction and its relation to the constructible topology
are recorded in Ferrand, Section 8.3, explicitly identifying it with
Grothendieck's $X^{\mathrm{cons}}$ in EGA I, 7.2.14
\cite{ROOT_Ferrand2003,ROOT_EGAI1971}.  On an affine it is
\[
 T(\operatorname{Spec}A)=\operatorname{Spec}
 \left(A[t_a:a\in A]/(a^2t_a-a,\,at_a^2-t_a:a\in A)\right).
\]
The two relations specify the unique pointwise inverse of $a$.
The universal property, residue fields and constructible topology follow
from Olivier's explicit construction
\cite[expos\'e 6]{ROOT_Olivier1968}\cite[Sections 1--3]{ROOT_Olivier1978}.
The morphism $t_X$ is affine, surjective and a monomorphism.

\begin{theorem}[Proved: comparison with finite piecewise morphisms]\label{C:absolute-flat-piecewise}
Let $k$ be algebraically closed, and let $X,Y$ be separated reduced
finite-type $k$-schemes.  There are natural bijections
\[
 \operatorname{Hom}_{\mathcal P_k}(X,Y)
 \simeq \operatorname{Hom}_{\operatorname{Sch}_k}(T(X),Y)
 \simeq \operatorname{Hom}_{\operatorname{Abs}_k}(T(X),T(Y)).
\]
Consequently $T$ gives a fully faithful realization of $\mathcal P_k$
as the full subcategory of $\operatorname{Abs}_k$ with objects $T(X)$
arising from separated finite-type reduced schemes (and their finite stratifications).
Its categorical product is
\[
 T(X)\mathbin{\times_{\operatorname{Abs}_k}}T(Y)
   \simeq T(X\times_kY)
   \simeq T\bigl(T(X)\times_kT(Y)\bigr).
\]
The product occurring inside the last $T$ is the ordinary scheme product.
It need not itself be absolutely flat.
\end{theorem}

\begin{proof}
We give the finite-partition limit argument, including the reduction
and finiteness points.  Let $\mathfrak P(X)$ be the directed system of
finite partitions of $X$ into affine locally closed reduced subschemes.
Such partitions exist by noetherian induction.  Common refinements
exist: intersect two partitions and, if necessary, subdivide the
intersections into finitely many affine locally closed reduced pieces.
Write $X_P=\coprod_{S\in P}S$.  Each $X_P$ is affine, of finite type,
and reduced.  For a refinement $P'$ of $P$ there is a natural morphism
$X_{P'}\to X_P$.  These transition morphisms are affine: the inverse
image of each affine component is a finite disjoint union of affine
components.  Let $L=\varprojlim_P X_P$; it is an affine scheme.

Its coordinate ring is absolutely flat.  A function $b$ in that ring
comes from some $\Gamma(X_P,\mathcal O)$, since the ring is a filtered
colimit.  Subdivide every component of $P$ into $D(b)$ and $V(b)_{
\mathrm{red}}$.  These are affine locally closed reduced subschemes,
so they form a refinement.  On the resulting disjoint union define
$b^{[-1]}$ to be $b^{-1}$ on $D(b)$ and zero on $V(b)_{
\mathrm{red}}$.  It satisfies
$b^2b^{[-1]}=b$ and $b(b^{[-1]})^2=b^{[-1]}$.  Thus every element of
the limit ring has a pointwise inverse.  Reduction of the closed
strata is essential here: no nilpotent thickening has been retained.

We next identify $L$ with $T(X)$ by the universal property.  If $Z$ is
absolutely flat and $f:Z\to X$ is a morphism, the inverse image of every
constructible subset of $X$ is clopen in $Z$.  This can be checked
locally on affine opens: every principal open in the spectrum of an
absolutely flat ring is clopen, and each constructible subset of the
noetherian scheme $X$ is a finite Boolean combination of such opens.
Moreover, a morphism from a reduced scheme whose image is contained
in a locally closed reduced subscheme factors through that subscheme;
on a closed affine piece, its defining ideal maps into all primes and
hence to zero.  It follows that $f$ factors uniquely through $X_P$
for every $P$, using the clopen inverse-image pieces of $Z$.  The
factorizations are compatible with refinement.  Hence $f$ factors
uniquely through $L$.  This proves $L\simeq T(X)$.

Since $Y$ is of finite presentation over $k$, the standard limit
theorem gives
\[
 \operatorname{Hom}_k(T(X),Y)
 =\varinjlim_P\operatorname{Hom}_k(X_P,Y).
\]
Here all hypotheses are present: the stages are quasi-compact and
quasi-separated, and the transition maps are affine.  One may apply
Stacks, Lemma 32.10.1, parts (2) and (3), tag 01ZM\cite[tag 01ZM]{BStacks}, to the section of
$Y\times_k T(X)\to T(X)$ corresponding to a morphism; both this
product and the identity source descend to each stage.  This proves
both existence of a finite stage and equality after a common
refinement.  A morphism $X_P\to Y$ is exactly a map regular on every
piece of $P$.  Two such maps agree as maps of $k$-points exactly when
their restrictions to a common refinement agree as scheme morphisms,
because the pieces are reduced and of finite type over an algebraically
closed field.  The colimit is therefore
$\operatorname{Hom}_{\mathcal P_k}(X,Y)$.  The second displayed
bijection is the defining adjunction for $T(Y)$.

For products, test against an arbitrary $Z\in\operatorname{Abs}_k$:
\[
\begin{aligned}
 \operatorname{Hom}(Z,T(X\times_kY))
  &\simeq\operatorname{Hom}(Z,X\times_kY)\\
  &\simeq\operatorname{Hom}(Z,T(X))\times
          \operatorname{Hom}(Z,T(Y)).
\end{aligned}
\]
This proves the internal product formula.  Applying the same argument
to the ordinary product $T(X)\times_kT(Y)$ proves the second formula.
The terminal object is $\operatorname{Spec}k$, so this fully faithful
realization preserves the finite products used to define group objects.
\end{proof}

\paragraph{Why the ordinary product cannot be substituted.}
The fields $k(t)$ and $k(s)$ are absolutely flat rings, but
$k(t)\otimes_k k(s)$ has both the zero prime and the nonzero prime
defined by $t-s$.  Indeed it is the localization of $k[t,s]$ obtained
by inverting all nonzero polynomials in $t$ alone and in $s$ alone;
none of these meets the prime $(t-s)$.  Thus this tensor product is
not absolutely flat.  The same issue already occurs for
$T(\mathbb A^1_k)\times_k T(\mathbb A^1_k)$: its coordinate ring
has the displayed tensor product as a quotient, obtained by taking
the generic residue-field quotient on both factors.  A quotient of
an absolutely flat ring would be absolutely flat, so the ordinary
product cannot be absolutely flat in this example.

\begin{corollary}[Proved: algebraization in the absolutely flat category]\label{C:absolute-flat-group-algebraization}
With the internal product just specified, group objects in the full
subcategory formed by the $T(X)$ for separated finite-type reduced $k$-schemes
are equivalent to algebraic $k$-groups.  In characteristic zero this
also realizes group objects with arbitrary constructible graphs.
\end{corollary}

\begin{proof}
Apply Theorem~\ref{C:absolute-flat-piecewise} to the multiplication,
inverse and unit morphisms and to their commutative diagrams.  Full
faithfulness makes the group identities equivalent.  Then apply
Theorem~\ref{C:piecewise-group-algebraization}.
\end{proof}

\paragraph{Exact historical scope.}
The reference to EGA I, 7.2.14 is an explicit published identification
of the universal absolutely flat construction with Grothendieck's
constructible-topology scheme.  The theorem above proves its comparison
with the stated \emph{finite-type, reduced, finite-partition} category.
It does not by itself identify the complete category intended by the
1964 subject, or extend the algebraization assertion to all absolutely
flat schemes, pro-affine schemes, or arbitrary sheaves.
